\section{第8章 向量代数与空间解析几何}
\subsection{向量及其线性运算}
\begin{exercise}{1}{基础·向量线性组合}
设 $\vec u=\vec a-\vec b+\vec c$、$\vec v=\vec a+3\vec b-2\vec c$。用 $\vec a,\vec b,\vec c$ 表示 $2\vec u-\vec v$。
\end{exercise}
\begin{exercise}{2}{基础·坐标运算}
设 $\vec a=(0,-1,1)$、$\vec b=(1,2,0)$、$\vec c=(2,0,1)$。求 $\vec a-\vec b+\vec c$ 及其模。
\end{exercise}
\begin{exercise}{3}{基础·三等分点}
已知 $A(0,0,0)$、$B(6,3,-3)$，求线段 $AB$ 的两个三等分点及线段长度。
\end{exercise}
\begin{exercise}{4}{基础·线段内分点}
已知 $A(1,0,2)$、$B(4,6,-1)$，点 $M$ 在线段 $AB$ 上且 $|AM|:|MB|=1:2$。求 $M$ 的坐标。
\end{exercise}
\begin{exercise}{5}{基础·向量平行条件}
设 $A(1,2,3)$、$B(3,5,t)$，$\vec a=(2,3,-2)$。求使 $\overrightarrow{AB}$ 与 $\vec a$ 平行的 $t$。
\end{exercise}
\begin{exercise}{6}{基础·平行向量坐标}
设 $\vec a=(1,2,3)$、$\vec b=(2,k,t)$。求使两向量平行的 $k,t$。
\end{exercise}
\begin{exercise}{7}{基础·单位向量}
设 $A(-1,2,2)$、$B(1,2,0)$。求与 $\overrightarrow{AB}$ 反向的单位向量。
\end{exercise}
\begin{exercise}{8}{基础·方向余弦}
设 $\vec a=(1,1,\sqrt2)$，求其模及三个方向余弦。
\end{exercise}
\begin{exercise}{9}{基础·标量投影}
向量 $\vec a$ 的模为 $6$，与某有向轴的夹角为 $\pi/3$。求它在该轴上的标量投影。
\end{exercise}
\begin{exercise}{10}{常规·模与正交条件}
设 $\vec a=(1,2,3)$、$\vec b=(2,-1,p)$。由 $|\vec a+\vec b|=|\vec a-\vec b|$ 求 $p$。提示：将两边平方。
\end{exercise}
\begin{exercise}{11}{基础·两点等距轨迹}
求与 $A(0,0,0)$、$B(2,0,0)$ 等距的空间动点 $M(x,y,z)$ 的轨迹方程。
\end{exercise}
\begin{exercise}{12}{常规·中位线向量证明}
三角形 $ABC$ 中 $M,N$ 分别为 $AB,AC$ 的中点。用 $\overrightarrow{MN}=\overrightarrow{AN}-\overrightarrow{AM}$ 证明 $\overrightarrow{MN}=\frac12\overrightarrow{BC}$。
\end{exercise}
\begin{exercise}{13}{基础·中点坐标变式}
求 $A(1,-2,3)$、$B(3,4,-1)$ 的中点坐标，并求原点到该中点的距离。
\end{exercise}
\begin{exercise}{14}{常规·线段交点小综合}
已知 $A(0,0,0)$、$B(2,2,0)$、$C(0,2,0)$、$D(2,0,0)$。写出直线 $AB,CD$ 的参数方程并求交点。
\end{exercise}
\subsection{内积、向量积与混合积}
\begin{exercise}{15}{基础·三种向量乘积}
设 $\vec a=(1,0,0)$、$\vec b=(0,2,0)$、$\vec c=(0,0,3)$。计算 $\vec a\cdot\vec b$、$\vec a\times\vec b$ 及 $(\vec a\times\vec b)\cdot\vec c$。
\end{exercise}
\begin{exercise}{16}{基础·混合积与体积}
设 $\vec a=(1,0,1)$、$\vec b=(0,1,1)$、$\vec c=(0,0,2)$。求三向量张成的平行六面体体积。
\end{exercise}
\begin{exercise}{17}{基础·向量积运算}
化简 $(\vec a+\vec b)\times(\vec a-\vec b)$，结果用 $\vec a\times\vec b$ 表示。
\end{exercise}
\begin{exercise}{18}{基础·内积与垂直}
设 $\vec a=(1,2,-1)$、$\vec b=(2,0,2)$，判断两向量是否垂直。
\end{exercise}
\begin{exercise}{19}{基础·向量夹角}
求 $\vec a=(1,1,0)$ 与 $\vec b=(1,0,1)$ 的夹角。
\end{exercise}
\begin{exercise}{20}{基础·三角形面积}
用向量积求 $A(0,0,0)$、$B(2,0,0)$、$C(0,3,0)$ 所成三角形的面积。
\end{exercise}
\begin{exercise}{21}{常规·面积与高}
设 $A(1,0,0)$、$B(0,1,0)$、$C(0,0,1)$。求三角形 $ABC$ 的面积及从 $A$ 到直线 $BC$ 的距离。
\end{exercise}
\begin{exercise}{22}{基础·共线与内积}
设 $\vec m=(2,-3,1)$。求与 $\vec m$ 共线且满足 $\vec r\cdot\vec m=14$ 的 $\vec r$。
\end{exercise}
\begin{exercise}{23}{基础·正交参数}
设 $\vec a=(1,2,0)$、$\vec b=(k,-1,3)$。求使 $\vec a\perp\vec b$ 的 $k$。
\end{exercise}
\begin{exercise}{24}{基础·内积与向量积关系}
设 $|\vec a|=3$、$|\vec b|=4$，两向量夹角为 $\pi/6$。求 $\vec a\cdot\vec b$ 与 $|\vec a\times\vec b|$。
\end{exercise}
\begin{exercise}{25}{常规·向量和的模}
设 $|\vec a|=2$、$|\vec b|=3$，夹角为 $2\pi/3$。求 $|\vec a+\vec b|$ 与 $|\vec a-\vec b|$。
\end{exercise}
\begin{exercise}{26}{基础·共面判定}
用混合积判断 $\vec a=(1,0,1)$、$\vec b=(0,1,1)$、$\vec c=(1,1,2)$ 是否共面。
\end{exercise}
\begin{exercise}{27}{基础·乘积恒等式核验}
设 $\vec a=(1,1,0)$、$\vec b=(0,1,1)$。计算并核验 $|\vec a\times\vec b|^2+(\vec a\cdot\vec b)^2=|\vec a|^2|\vec b|^2$。
\end{exercise}
\begin{exercise}{28}{常规·余弦定理}
利用 $\overrightarrow{BC}=\overrightarrow{AC}-\overrightarrow{AB}$ 及向量模的平方，证明三角形的余弦定理 $a^2=b^2+c^2-2bc\cos A$。
\end{exercise}
\begin{exercise}{29}{基础·两向量张成面积}
设 $\vec a=(1,0,1)$、$\vec b=(0,2,0)$。求以它们为邻边的平行四边形面积。
\end{exercise}
\begin{exercise}{30}{基础·四面体体积变式}
求以 $O(0,0,0)$、$A(2,0,0)$、$B(0,3,0)$、$C(0,0,4)$ 为顶点的四面体体积。
\end{exercise}
\begin{exercise}{31}{常规·混合积参数小综合}
设 $\vec a=(1,1,0)$、$\vec b=(0,1,1)$、$\vec c=(1,t,1)$。求使三向量共面的 $t$，并在此时将 $\vec c$ 表示为 $\vec a,\vec b$ 的线性组合。
\end{exercise}
\subsection{平面方程}
\begin{exercise}{32}{基础·点法式平面}
求过 $A(1,2,3)$、法向量为 $(2,1,-1)$ 的平面方程。
\end{exercise}
\begin{exercise}{33}{基础·三点确定平面}
求过 $A(1,0,0)$、$B(0,2,0)$、$C(0,0,3)$ 的平面方程。
\end{exercise}
\begin{exercise}{34}{基础·平行坐标轴的平面}
求平行于 $y$ 轴且过 $A(1,0,1)$、$B(2,-1,3)$ 的平面方程。
\end{exercise}
\begin{exercise}{35}{基础·特殊平面}
写出平面 $3y-1=0$ 的一个法向量，并指出它平行于哪个坐标面。
\end{exercise}
\begin{exercise}{36}{常规·垂直平面}
求过 $A(0,0,0)$、$B(1,1,0)$ 且垂直于 $x+y+z=1$ 的平面方程。
\end{exercise}
\begin{exercise}{37}{基础·双重垂直条件}
求过原点且同时垂直于平面 $x=1$、$y=2$ 的平面方程。
\end{exercise}
\begin{exercise}{38}{基础·垂直平分面}
求 $A(1,2,3)$、$B(3,2,1)$ 的垂直平分面。
\end{exercise}
\begin{exercise}{39}{基础·两方向确定平面}
求过原点且平行于向量 $(1,1,0)$、$(0,1,1)$ 的平面方程。
\end{exercise}
\begin{exercise}{40}{基础·截距式}
求在 $x,y,z$ 轴上的截距分别为 $2,3,6$ 的平面方程，并求原点到该平面的距离。
\end{exercise}
\begin{exercise}{41}{基础·平面夹角}
求平面 $x+y=0$ 与 $x-y=1$ 的夹角。
\end{exercise}
\begin{exercise}{42}{基础·坐标面夹角}
求平面 $x+y+\sqrt2z=1$ 与 $xOy$ 面的夹角。
\end{exercise}
\begin{exercise}{43}{基础·过点的平行平面}
求过 $P(1,0,0)$ 且平行于 $2x+y-z=3$ 的平面方程。
\end{exercise}
\begin{exercise}{44}{基础·平面位置关系}
判断 $x+2y-z=1$ 与 $2x+4y-2z=3$ 的位置关系；若平行，求两平面距离。
\end{exercise}
\begin{exercise}{45}{基础·平面交线}
求平面 $x+y=2$ 与 $x=1$ 的交线参数方程。
\end{exercise}
\begin{exercise}{46}{基础·平面法向量}
求平面 $2x+y-2z=5$ 与三个坐标面的夹角余弦。
\end{exercise}
\begin{exercise}{47}{常规·点面距离与垂足}
求 $P(2,0,0)$ 到平面 $x+y+z=1$ 的距离及垂足坐标。
\end{exercise}
\begin{exercise}{48}{基础·等距平行平面}
求与平面 $z=1$ 平行且相距 $2$ 的两个平面方程。
\end{exercise}
\begin{exercise}{49}{基础·简单切平面}
求曲面 $z=x^2+y^2$ 在 $(1,1,2)$ 处的切平面。
\end{exercise}
\begin{exercise}{50}{基础·截距与体积}
求平面 $x/2+y/4-z/3=1$ 与三个坐标面围成四面体的体积。
\end{exercise}
\begin{exercise}{51}{基础·平行平面距离}
求 $3x+6y-2z-7=0$ 与 $3x+6y-2z+14=0$ 的距离。
\end{exercise}
\begin{exercise}{52}{常规·球面与平面相切}
球心为 $O(1,1,1)$ 的球面与平面 $x+y+z=0$ 相切。求球面半径与切点。
\end{exercise}
\begin{exercise}{53}{常规·过直线的平面}
求包含直线 $(x,y,z)=(1,0,0)+t(1,1,0)$ 且平行于向量 $(0,0,1)$ 的平面方程。
\end{exercise}
\begin{exercise}{54}{常规·角平分面}
求平面 $x-y=0$ 与 $x+y=2$ 的两个角平分面，提示：令点到两平面的距离相等。
\end{exercise}
\begin{exercise}{55}{基础·平面束变式}
求过直线 $x=y=0$ 且过点 $P(1,2,3)$ 的平面方程。
\end{exercise}
\begin{exercise}{56}{常规·平面与截体小综合}
平面过 $P(1,1,1)$，且在三个正半轴上的截距相等。求平面方程及它与三个坐标面围成四面体的体积。
\end{exercise}
\subsection{空间直线}
\begin{exercise}{57}{基础·两点式直线}
求过 $A(1,2,3)$、$B(3,2,1)$ 的直线参数方程，注意方向向量的 $y$ 分量为零。
\end{exercise}
\begin{exercise}{58}{基础·平面法线}
求过 $P(1,0,2)$ 且垂直于 $x-2y+z=3$ 的直线参数方程。
\end{exercise}
\begin{exercise}{59}{基础·一般式转参数式}
将直线 $x+y=1$、$z=2$ 写成参数方程。
\end{exercise}
\begin{exercise}{60}{基础·公共平行方向}
求过原点且方向同时平行于平面 $x+y=1$ 和 $y+z=2$ 的直线参数方程。
\end{exercise}
\begin{exercise}{61}{基础·公共垂直方向}
求过 $P(1,1,1)$ 且方向同时垂直于 $(1,1,0)$ 与 $(0,1,1)$ 的直线参数方程。
\end{exercise}
\begin{exercise}{62}{基础·直线方向角}
求直线 $(x,y,z)=(1,2,0)+t(1,1,\sqrt2)$ 与 $z$ 轴的夹角。
\end{exercise}
\begin{exercise}{63}{常规·点到直线的垂足}
求 $P(1,0,0)$ 在直线 $(x,y,z)=t(1,1,0)$ 上的垂足及点线距离。
\end{exercise}
\begin{exercise}{64}{基础·点线距离}
求 $P(1,2,3)$ 到 $z$ 轴的距离及垂足坐标。
\end{exercise}
\begin{exercise}{65}{基础·两直线夹角}
求方向向量分别为 $(1,0,1)$、$(0,1,1)$ 的两条直线的夹角。
\end{exercise}
\begin{exercise}{66}{常规·一般式直线方向}
求直线 $x+y+z=0$、$x-y=0$ 的一个方向向量，并求该直线与 $z$ 轴的夹角。
\end{exercise}
\begin{exercise}{67}{基础·线面夹角}
求直线 $(x,y,z)=t(1,0,1)$ 与平面 $z=0$ 的夹角。
\end{exercise}
\begin{exercise}{68}{基础·方向正交参数}
两直线方向向量为 $(1,t,2)$ 和 $(4,2,-1)$。求使两方向垂直的 $t$。
\end{exercise}
\begin{exercise}{69}{基础·线面交点}
求直线 $(x,y,z)=(1,0,0)+t(1,1,1)$ 与平面 $x+y+z=4$ 的交点。
\end{exercise}
\begin{exercise}{70}{基础·直线包含于平面}
判断直线 $(x,y,z)=(1,0,0)+t(1,-1,0)$ 是否位于平面 $x+y=1$ 内。
\end{exercise}
\begin{exercise}{71}{基础·代入参数求交点}
求直线 $(x,y,z)=(0,1,2)+t(2,-1,1)$ 与平面 $x+y+z=6$ 的交点。
\end{exercise}
\begin{exercise}{72}{基础·线面平行判定}
判断直线 $(x,y,z)=(0,0,1)+t(1,-1,0)$ 与平面 $x+y+z=0$ 的位置关系。
\end{exercise}
\begin{exercise}{73}{常规·平面内指定方向直线}
求经过原点、位于平面 $x+y+z=0$ 内且方向垂直于 $(1,-1,0)$ 的直线方程。
\end{exercise}
\begin{exercise}{74}{基础·线面距离}
求直线 $(x,y,z)=(0,0,2)+t(1,1,0)$ 到平面 $z=0$ 的距离。
\end{exercise}
\begin{exercise}{75}{基础·平行线面距离}
求直线 $(x,y,z)=(1,0,0)+t(1,-1,0)$ 与平面 $x+y+z=3$ 的距离。
\end{exercise}
\begin{exercise}{76}{常规·异面直线距离}
求 $L_1:(x,y,z)=(t,0,0)$ 与 $L_2:(x,y,z)=(0,s,2)$ 的距离，并指出一对垂足。
\end{exercise}
\begin{exercise}{77}{常规·点的平面投影}
求 $P(1,2,3)$ 在平面 $x+y+z=0$ 上的投影点及关于该平面的对称点。
\end{exercise}
\begin{exercise}{78}{常规·点的直线投影}
求 $P(1,0,1)$ 在直线 $(x,y,z)=t(1,1,1)$ 上的投影点及点线距离。
\end{exercise}
\begin{exercise}{79}{基础·平行与包含}
判断直线 $(x,y,z)=(1,0,0)+t(1,-1,0)$ 与平面 $x+y=1$、$x+y=2$ 的位置关系。
\end{exercise}
\begin{exercise}{80}{常规·简单公垂线}
设 $L_1:(x,y,z)=(t,0,0)$、$L_2:(x,y,z)=(0,1,s)$。求两直线的公垂线与距离。
\end{exercise}
\begin{exercise}{81}{基础·直线夹角变式}
求直线 $(x,y,z)=t(1,1,0)$ 与 $(x,y,z)=(0,0,1)+s(1,-1,0)$ 的夹角。
\end{exercise}
\begin{exercise}{82}{常规·位置距离小综合}
设 $L_1:(x,y,z)=(t,0,0)$、$L_2:(x,y,z)=(0,1,2)+s(0,1,1)$。判断两直线是否异面，并求它们的距离。
\end{exercise}
\subsection{曲面及其方程}
\begin{exercise}{83}{基础·配方识别球面}
求球面 $x^2+y^2+z^2-6x+8y+2z+10=0$ 的球心与半径。
\end{exercise}
\begin{exercise}{84}{基础·直径端点求球面}
球面一条直径的端点为 $A(2,-3,5)$、$B(4,1,-3)$。求球面方程。
\end{exercise}
\begin{exercise}{85}{常规·四点求球面}
求过 $O(0,0,0)$、$A(2,0,0)$、$B(0,2,0)$、$C(0,0,2)$ 的球面方程。
\end{exercise}
\begin{exercise}{86}{基础·标准圆柱面}
求以 $z$ 轴为轴线、半径为 $2$ 的圆柱面方程，并写出它与 $z=1$ 的截线方程。
\end{exercise}
\begin{exercise}{87}{基础·二次曲面识别}
识别 $x^2+y^2-z^2=1$、$x^2-y^2-z^2=1$ 的曲面类型，并指出它们的轴向。
\end{exercise}
\begin{exercise}{88}{基础·旋转球面}
将 $xOz$ 面内的圆 $x^2+z^2=9$ 绕 $z$ 轴旋转一周，求所得曲面方程。
\end{exercise}
\begin{exercise}{89}{基础·旋转双曲面}
将 $xOy$ 面内的双曲线 $x^2-y^2=1$ 绕 $y$ 轴旋转，求所得曲面方程并识别类型。
\end{exercise}
\begin{exercise}{90}{基础·缺失变量与柱面}
说明空间中的 $y^2=x$、$x^2+y^2=1$ 各是什么柱面，并指出母线方向。
\end{exercise}
\begin{exercise}{91}{基础·标准二次曲面}
将 $x^2+2y^2+3z^2=6$ 化为标准形式，指出曲面类型及三个半轴长。
\end{exercise}
\begin{exercise}{92}{常规·球面相切}
求半径为 $2$ 且与平面 $z=0$ 相切于原点的两个球面方程。
\end{exercise}
\begin{exercise}{93}{常规·旋转椭球}
将 $xOy$ 面内的椭圆 $x^2/4+y^2/3=1$ 绕 $x$ 轴旋转。求椭球面方程，并求它与平面 $x=1$ 的截线半径。
\end{exercise}
\begin{exercise}{94}{基础·抛物面变式}
识别 $z=x^2+2y^2$ 的曲面类型，指出顶点，并写出 $z=2$ 的截线方程。
\end{exercise}
\begin{exercise}{95}{常规·交叉项小综合}
利用 $u=(x+z)/\sqrt2$、$v=(x-z)/\sqrt2$，将 $x^2+y^2+z^2+xz=1$ 化为标准形式，识别曲面类型。
\end{exercise}
